\documentclass[12pt]{article}
\usepackage[utf8]{inputenc}
\usepackage[T1]{fontenc}
\usepackage{hyperref}
\title{Kachari Language Resources for K-12 Curriculum Development}
\author{Compiled by Claude Sonnet 5.5 (claude-sonnet-5-5) for Andr\'as Kornai}
\date{}
\begin{document}
\maketitle

\section{Introduction}
Kachari (ISO 639-3 xac; Glottocode kach1279) is a Sino-Tibetan language of Assam, India. The name ``Kachari'' is a generic label applied to several ethnic groups in Assam who speak Tibeto-Burman languages or claim a shared ancestry. Because of this, it is easy to confuse the language with Dimasa (Dimasa Kachari) or with the Sonowal Kachari, who are reported to use Assamese mainly. The search did not establish which resources belong specifically to ISO code xac, so none of the materials below should be assumed to be in the xac language.

\section{Primary Educational Resources}
No primers or textbooks verified as belonging to xac were found. Search results mention that Dimasa, a separate Boro-Garo language, is taught as a Mother-tongue/Modern Indian Language subject in Dima Hasao, and that Dimasa textbooks reportedly exist. That is Dimasa, not necessarily xac, and no URL was found for those textbooks.

\section{Linguistic and Reference Materials}
\begin{itemize}
\item Wikipedia on the Dimasa Kachari language (a related but separate language; treat with caution): \url{https://en.wikipedia.org/wiki/Dimasa_Kachari_language}
\item Wikipedia on the Sonowal Kachari people: \url{https://en.wikipedia.org/wiki/Sonowal_Kachari_people}
\item Endangered Languages Project entry surfaced for Kachari: \url{https://endangeredlanguages.com/node/91704}
\end{itemize}

\section{Religious and Community Resources}
Joshua Project lists the Sonwal Kachari and describes language and scripture status for that group: \url{https://joshuaproject.net/index.php/people_groups/17296/IN}. This is outreach material, and its link to xac is unverified.

\section{Digital and Audio Resources}
No audio, video, corpus or software resources were found for xac.

\section{Curriculum Development Context}
Little verified material exists for this code in the search results. A curriculum would first need to identify which community and variety xac refers to, and then work with community members and Assam educational authorities. Resources for neighbouring languages (for example Dimasa) should not be presented as Kachari without confirmation.

\end{document}